\begin{frame}{Síntese dos resultados e limites da comparação}
  \small
  {\centering
    \setlength{\tabcolsep}{5pt}
    \renewcommand{\arraystretch}{0.95}
    \begin{tabular}{lrrrrrr}
      \toprule
      & \multicolumn{3}{c}{Frequência de corte (Hz)} & \multicolumn{3}{c}{Ganho de passagem (V/V)} \\
      \cmidrule(lr){2-4}\cmidrule(lr){5-7}
      Circuito & Teoria & Medido & Erro & Teoria & Medido & Erro \\
      \midrule
      Sem carga & $70{,}20$ & $67$ & $4{,}56\%$ & $1{,}000$ & $0{,}971$ & $2{,}92\%$ \\
      Carga direta & $109{,}58$ & $121$ & $10{,}42\%$ & $0{,}595$ & $0{,}404$ & $32{,}15\%$ \\
      Buffer & $68{,}83$ & $69$ & $0{,}25\%$ & $1{,}000$ & $0{,}960$ & $4{,}00\%$ \\
      \bottomrule
    \end{tabular}\par
  }

  \centering\scriptsize Erros relativos à teoria com componentes medidos, calculados antes do arredondamento.
  \par\small
  \begin{columns}[T,onlytextwidth]
    \begin{column}{0.48\textwidth}
      \begin{block}{Resultados sustentados}
        \begin{itemize}
          \item A carga direta reduz o ganho e desloca o corte.
          \item O buffer aproxima a resposta daquela sem carga.
          \item A tensão de entrada medida permite comparar toda a varredura.
        \end{itemize}
      \end{block}
    \end{column}
    \begin{column}{0.48\textwidth}
      \begin{alertblock}{Limites da interpretação}
        \begin{itemize}
          \item O ganho do circuito 2 diverge do modelo em $32{,}15\%$.
          \item Convém conferir componentes, conexões e referência de entrada.
          \item Os dados não isolam uma causa para esse desvio.
        \end{itemize}
      \end{alertblock}
    \end{column}
  \end{columns}
\end{frame}
